\documentclass[11pt]{ctexart}
\input{D:/aaa_D_usual/DSH/111学习助手/.dsh/skills/lecture-deepdive/assets/theme-violet}
\begin{document}
\doctitle{杨氏双缝：从光程差到云母片}%
{第五章 \S5.5--5.7 \ 双缝干涉与分波前装置 · 精讲讲义}%
{依据：\texttt{第五章 光的干涉.pdf}（83 页扫描版，逐页读图）p13--p28\ \textbar\ 2026-10-27}
\callout{\textbf{这份讲义解决什么。}你手上那张手写稿（云母片盖住一条缝、算出 $7.6\,\mu\mathrm{m}$）对应的不是习题册，而是课件本身：\textbf{p20 的 Ex.\,5.2，页脚印着 HOME WORK TWO}。整份作业只有这一问。
本讲义把 p13--p28 这 16 页串成一条能独立读的线：\S1 讲清"普通光凭什么也能干涉"；\S2 把光程与相位差接上；\S3 补完双缝的几何近似（课件在这里跳了一步，是全班最容易写错的地方）；\S4、\S5 逐题全解（\S5 就是你的作业，给了三条互相独立的路线）；\S6--\S8 讲对比度、空间相干性和三种分波前装置；\S11 把课件每一处跳步连页码带原话列出来；\S12 给三题自检。
讲义里凡标 \textbf{[补]} 的地方，都是课件没有写、由本讲义补出的完整推导；凡标 \textbf{[课件原话]} 的地方，都是直接引用那 16 页的句子。}

\[
\vec{E}(P,t)=\vec{E}_1(P,t)+\vec{E}_2(P,t).
\]

\end{document}
